\documentclass[11pt,a4paper]{article}
\usepackage[margin=25mm]{geometry}
\usepackage{amsmath,amssymb}
\setlength{\parindent}{0pt}
\setlength{\parskip}{0.6em}
\setlength{\emergencystretch}{2em}
\begin{document}
\begin{center}
  {\LARGE\bfseries G --- Gerrymandering}\\[0.5em]
  {\small Two pointers, matrices, DP, sliding window}
\end{center}

\subsection*{1. The main idea}

We need to count the subregions whose corruptness is at least $c$.
The main idea is that corruptness cannot go down when we extend a
subregion. If we can fix a larger subregion, then we have also fixed
every smaller subregion inside it, using at most that many changes.

So the problem boils down to this: for every left endpoint $l$, find the
first right endpoint $r$ where we need at least $c$ changes. Every right
endpoint from $r$ to $n$ also works, so we add $n-r+1$ to the answer.
The columns are numbered from $1$ to $n$.

We can do this with two pointers since moving $l$ to the right cannot
make the first working $r$ move left. Keep extending the window until
we need at least $c$ changes, count the intervals, then remove the leftmost
column and repeat. If we reach the end and still have fewer than $c$
changes, we can stop. Now we just need a data structure to tell us the
minimum number of changes for the current window.

\subsection*{2. The transition table}
Since there are only three rows, there are only $3\cdot2\cdot2=12$
valid ways to color one column. The top and bottom can have the same
color. We only care about adjacent cells. These 12 configurations are
our DP states. Two states can go next to each other if they have
different colors in all three rows.

Let $w_i(b)$ be the number of changes needed to turn column $i$ into
state $b$. For a fixed interval, let $dp_i(b)$ be the minimum cost so far
if column $i$ ends up in state $b$. The transition is just
\[
  dp_i(b)=w_i(b)+\min_{a\text{ compatible with }b}dp_{i-1}(a).
\]
For the first column, we start with $dp_l(b)=w_l(b)$.

\subsection*{3. Turning the DP into matrices}

We can put the transition into a $12\times12$ matrix:
\[
  A_i[a,b]=\begin{cases}
    w_i(b),&a\text{ and }b\text{ are compatible},\\
    \infty,&\text{otherwise}.
  \end{cases}
\]
To combine two transitions, try every state $k$ between them and take
the cheapest option:
\[
  (A\otimes B)[a,b]=\min_k\bigl(A[a,k]+B[k,b]\bigr).
\]
This is called \emph{min-plus matrix multiplication}. Combining all the
matrices from $l$ to $r$ gives us the DP for that window, and the smallest
entry is the answer for its corruptness.

The base case is $dp_l(b)=w_l(b)$. The first column has no restriction
from a previous column.

\subsection*{4. The segment tree was almost enough}

\begin{quote}
\small\itshape
We had the segtree solution during the contest, but it was too slow.
I even asked my teammate if there was a trick where we only stored
the nodes inside the window. Turns out the thing we needed was just
two stacks. Mogged by two stacks.
\end{quote}

A segment tree can maintain this matrix product, but it adds a
$\log n$ factor. Here we only add columns on the right and remove them
on the left. We never change anything in the middle. Think of the
window as a deque, where each move changes just one end:

\begin{center}
\renewcommand{\arraystretch}{1.7}
\begin{tabular}{rc}
Current window &
  $\boxed{A_l}\;\boxed{A_{l+1}}\;\boxed{\cdots}\;\boxed{A_r}$ \\
& $\Big\downarrow\quad\text{push one column on the right}$ \\
Extend $r$ &
  $\boxed{A_l}\;\boxed{A_{l+1}}\;\boxed{\cdots}\;\boxed{A_r}\;\boxed{A_{r+1}}$ \\
& $\Big\downarrow\quad\text{pop one column from the left}$ \\
Advance $l$ &
  $\boxed{A_{l+1}}\;\boxed{\cdots}\;\boxed{A_r}\;\boxed{A_{r+1}}$
\end{tabular}
\end{center}
We may push several times before popping, but each individual move
adds or removes only one column. We only use the deque as a queue:
push at the back, pop at the front.

This is where \emph{SWAG (sliding window aggregation)} comes in. It lets
us maintain the combined value of a queue using two stacks. The only
thing we need is an associative operation, meaning we can change the
grouping without changing the answer. Our matrix multiplication has
this property. SWAG keeps the matrices in their original order and
supports each window move using an amortized constant number of merges.
So it gives us exactly what the two-pointer approach needs, without
the extra $\log n$ factor.

\subsection*{5. Complexity}

Both pointers move at most $n$ times. With SWAG, this means we only
do $O(n)$ matrix multiplications overall. Each one takes $O(12^3)$ time,
so the time complexity is $O(12^3n)$ and requires $O(12^2n)$ space.

\end{document}
